% !TeX root = ../../protocolo.tex


In Fig. \sfref{fig:extinction}{a)} the experimental extinction efficiencies (scale on the left axis for orange lines) at normal incidence, as a function of the incident wavelength for the metasurface characterized by  $\mu = (5.78 \pm 0.06)$~nm in the top frame and by  $\mu = (7.30 \pm 0.05)$~nm in the bottom frame. These measurements are contrasted with the theoretical extinction efficiencies (scale on the right axis) of an isolated AuNP, calculated through standard Mie Theory, with radii $a = 5.78$~nm (red lines) and $a = 7.30$~nm (green lines). The calculations consider a matrix of either air (long dashed lines) with $n_\text{m} = 1$ or glass (short dashed lines) with $n_\text{m} = 1.5$. The refractive index used for the AuNP $n_\text{NP}(\omega)$ was obtained from the experimental data reported by Johnson and Christy \cite{johnson_optical_1972}, modified with a size correction \cite{mendoza_herrera_determination_2014} to the Drude contribution of the complex-valued dielectric function $\varepsilon_\text{NP}(\omega) = n_\text{NP}^2(\omega)$. 

\begin{figure}[b!]\centering
	\includegraphics[width = \textwidth]{non-int/Fig-Extinction.pdf}
	\caption{%
		\textbf{a)} Experimental extinction efficiency of the metasurfaces (continuous orange lines) when illuminated at normal incidence, and extinction efficiency of a single AuNP of radius $a = 5.78$~nm (red lines) and $a = 7.30$~nm (green lines) embedded in air (long dashed lines) and in glass (short dashed lines), as given by Mie Theory, as a function of the incident light wavelength. \textbf{b)} Extinction efficiencies as a function of the incidence wavelength, at normal incidence, of a single AuNP embedded in an infinite air matrix (continuous lines) calculated by means of Mie Theory, and in a semiinfinite air matrix supported on a semiinfinte glass substrate (markers) using FEM simulations. The red lines correspond to the calculations of a AuNP with radius $a = 5.78$~nm and the green lines to $a=7.30$~nm. The magenta markers show the excitation of the LSPR for each case.}
	\label{fig:extinction}
\end{figure}

From the experimental extinction efficiencies,  a resonance is observed at a wavelength of approximately 542nm (546nm) for the metasurfaces characterized by a nominal mean radius of $\mu=5.78$~nm ($\mu=7.30$~nm), shown in the upper (lower) frame. The theoretical extinction efficiencies for isolated AuNPs are found to be maximized at around $510$~nm when $n_m=1$ and at approximately $533$~nm when $n_=1.5$, for both chosen values of the nanoparticle radius $a$. These wavelengths correspond to the excitation of the electric dipolar localized surface plasmon resonance (LSPR). When the extinction spectra of the metasurfaces in Fig. \sfref{fig:optical}{a)} are compared to those of the corresponding isolated AuNPs, a redshift in the resonance wavelength of the metasurfaces is evident. If the AuNPs in the metasurfaces were treated as independent scatterers, a redshift of approximately $36$~nm would be expected when the nanoparticles are embedded in air. However, such a shift is not anticipated, even when the interaction between the particles and the substrate is taken into account.

To verify whether the spectral position of the LSPR observed in the experimental extinction spectra is influenced solely by the interaction between isolated AuNPs and the substrate, FEM calculations were performed. In Fig. \sfref{fig:extinction}{b)}, extinction efficiencies (solid lines) were calculated using standard Mie Theory for isolated AuNPs embedded in air considering radii of $a=5.78$~nm (red) and $a=7.30$~m (green) alongside the  FEM simulations (red and green markers connected by thin lines), where the AuNPs were placed on a glass substrate ($n_s=1.5$) and illuminated at normal incidence; magenta markers indicate the extinction efficiencies at the LSPR excitation wavelengths. It is observed that the LSPR appears at $\sim510$ nm in the Mie Theory calculations for both radii, while the presence of the substrate causes a redshift of only $\sim3$ nm. These results indicate that the resonances seen in the experimental extinction spectra of the metasurfaces in Fig. \sfref{fig:extinction}{a)} do not correspond to single-particle responses, but rather to collective excitations of the AuNPs forming the monolayer. To confirm that the spectral position of the observed resonances in the metasurfaces are a consequence of the multiple scattering among the AuNPs in the array, measurements of the reflectivity of the metasurface under ATR condition were performed, so the AuNPs of the metasurface were illuminated with an evanescent wave enhancing its optical properties.

In Fig. \ref{fig:Rspectra}, reflectivity spectra of synthesized metasurfaces were measured in an ATR configuration as a function of incident wavelength. Measurements were conducted using $s$-polarized (top row) and $p$-polarized (bottom row) light at angles of incidence $\theta_i=64.9^\circ$ and $\theta_i=72.9^\circ$, both exceeding the critical angle of a glass-air interface. In addition, theoretical reflectivity curves (opaque lines) were calculated with the CSM with nominal parameters from Fig. \ref{fig:morpho}. Moreover, a fitting of the CSM to the experimental data was performed to determine the best-fitting mean radius $a_\text{Fit}^\text{CSM}$ and best-fitting surface coverage $\Theta_\text{Fit}^\text{CSM}$ by minimizing the squared error under two constraints: a fixed surface coverage $\Theta_\text{Fit}^\text{CSM}$ for both polarizations per sample, and $\Theta_\text{Fit}^\text{CSM}0<0.30$, to ensure model validity. The resulting fit parameters are shown in each panel of Fig. \ref{fig:Rspectra}.

\begin{figure}[h]\centering\hspace*{-.5em}%
	\includegraphics[width=\textwidth]{non-int/Fig-Fit.pdf}
	\caption{Reflectivity at ATR configuration of a AuNPs metasurface as a function of the incident wavelength considering an angle of incidence $\theta_i = 64.9^\circ$ (left column) and $\theta_i = 72.9^\circ$  (right column), for $s$ (upper row) and $p$ (lower row) polarization state. The metasurface is embedded in a water matrix and supported onto a glass substrate with refractive indices $n_\text{m} = 1.33$  and $n_\text{s} = 1.5 $, respectively. The experimental  reflectivity (continuous translucent lines) of the metasurface is shown for the fabricated samples characterized by the mean radius $\mu$ of their size distribution and by its surface coverage $\Theta$. The red lines correspond to the metasurface shown in Fig. \ref{fig:morpho}\textbf{a)} while the green lines to the sample shown in Fig. \ref{fig:morpho}\textbf{b)}. The theoretical curves (opaque continuous and dashed lines) were calculated by the CSM considering the nominal values of the metasurface, and a mean radius $a^\text{CSM}$ for the size distribution and a surface coverage $\Theta^\text{CSM}$, respectively; these last values were obtained by fitting  the CSM to the experimental data through the descent gradient method.}
	\label{fig:Rspectra}
\end{figure}

By compairing the theoretical experimental results with the calculations permorfed with the CSM employing the nominal valuse of the metasurfaces, a qualitative agreement related to surface coverage was observed: higher surface coverage resulted in greater light extinction in the visible reflection spectra. Also, a dependence of the spectral shift magnitude on the angle of incidence was also captured, consistent with experimental results for $p$-polarized light. Specifically, theoretical redshifts between samples were estimated to be approximately $3$~nm and $4$~nm for incidence angles of $\theta_i=64.9^\circ$ and $\theta_i=72.9^\circ$, respectively. In contrast, no spectral shift between samples was predicted for $s$-polarized light at either angle, with all resonances appearing near $525$~nm. Additionally, the overall reflectivity predicted by the CSM using nominal values was higher than the experimental measurements for ss polarization, and lower for pp polarization.		

The values of the best-fitting parameters show the following trends for all considered cases. First, $a^\text{CSM}_\text{Fit}$ was found to be smaller than the nominal mean radius $\mu$: For $s$-polarization, $a^\text{CSM}_\text{Fit}$ increased with the angle of incidence, while for $p$-polarization, it decreased. The fitted values satisfied the small particle approximation, consistent with the lack of spectral shift between samples observed in Fig.~\ref{fig:Rspectra}. The fitted surface coverage $\Theta^\text{CSM}_\text{Fit}$ also showed polarization dependence. Values greater than the nominal characterization were obtained for $s$-polarization, while lower values were obtained for $p$-polarization. However, all fitted parameters remained close to those obtained from morphological characterization, suggesting reasonable consistency between models and experimental data.

Reflectivities calculated using the best-fitting parameters (dashed lines in Fig.~\ref{fig:Rspectra}) showed a reduction in $s$-polarized reflectivity relative to predictions based on nominal values, though both exceeded the experimental data. For $p$-polarization, the fitted model showed improved agreement with experiments across the visible spectrum, particularly for the sample with $\mu = 5.78$~nm at $\theta_i = 72.9^\circ$. These results indicate that, while the CSM with fitted parameters improves the match with experimental reflectivity, it does not fully capture the optical response of the metasurfaces. Notably, polarization-dependent behavior observed in experiments is not reproduced by the CSM or FEM simulations. This discrepancy suggests the presence of anisotropy in the system, potentially caused by partial embedding of the nanoparticles---an effect previously reported for thermally annealed AuNPs~\cite{meng2015anisotropic}.

To evaluate this effect, FEM simulations were performed for a single partially embedded AuNP with radius $a$, positioned at a height $h$ relative to the air-glass interface and illuminated at normal incidence. This configuration, shown in Fig.~\sfref{fig:roots}{a)}, uses $r$ to represent the apparent radius seen in SEM micrographs. To take into account a partial embedding of the AuNPs, it was considered an equivalent AuNP to be employed in the CSM with volume $V_\text{eq}$ and radius $a_\text{Fit}^\text{CSM}$, which are related to the fraction of the AuNPs above the substrate $	V_\text{above}^\text{semi-sphere}$ as
%
\begin{align}
	V_\text{above}^\text{semi-sphere} = \frac43 \pi a^3 \left[\frac12 - \frac34\left(\frac{h}{a}\right) + \frac14\left(\frac{h}{a}\right)^3\right] = \frac43 \pi (a_\text{Fit}^\text{CSM})^3 = V_\text{eq},
	\label{eq:volume}
\end{align}
%
where $h$ is the distance from the center of the AuNP to the glass-air interface; $h=a$ corresponds to a sphere totally embedded in the substrate and $h=-a$ to one perfectly supported on it.  The partial embedding of the metasurface can be calculated, in average, by solving for $h$ in Eq. \eqref{eq:volume}, which is equivalent  to calculate the roots of the function $\eta$, given by
%
\begin{align}
	\eta\left(\frac{h}{a}\right) = \frac12 - \frac34\left(\frac{h}{a}\right) + 
	\frac14\left(\frac{h}{a}\right)^3 -  \left(\frac{a_\text{Fit}^\text{CSM}}{a}\right)^3
	=0.
\end{align}
%	
In particular, in Fig. \sfref{fig:roots}{b)} the function $\eta(h/a)$ is plotted considering the average of $a_\text{Fit}^\text{CSM}$ for each metasurface characterized by the most probable radius $a$ of their AuNPs. The respective values of such parameters are presented in Fig. \sfref{fig:roots}{b)} and its roots are indicated with the blue markers. Since in both considered cases the roots are positive, it can be concluded that the most of the volume of the AuNPs of the metasurface are, in average, inside the substrate.

\begin{figure}[h!]\centering
	\includegraphics[scale = 1]{non-int/Fig-MieEmb.pdf}
	\caption{\textbf{a)} Schematics of a AuNP of radius $a$ partially embedded in a substrate, illuminated by an electromagnetic plane wave with an incident wave vector $\mathbf{k}_i$ perpendicular to the matrix-substrate interface. The distance between the center of the AuNP and the interface is defined as $h$ while $r$ is the aparent radius of the AuNP as seen from above. \textbf{b)} Plot of function $\eta(h/a)$ defined in Eq. \eqref{eq:volume} for a metasuface characterized for the most probable radius in its size distribution $a=5.78$ nm and the average of the best fittig radius to the CSM $\langle a_\text{Fit}^\text{CSM}\rangle = 3.3$ nm (red lines) and for a metasurfae with $a=7.30$ nm and $\langle a_\text{Fit}^\text{CSM}\rangle = 4.5$ nm (green lines). The roots of $\eta(h/a)$ are signilized with the blue markers for each case.}
	\label{fig:roots}
\end{figure}

Fig. \ref{fig:embedding} presents the extinction efficiencies of single AuNPs with radii  $a=5.78$ nm [Fig. \sfref{fig:embedding}{a)}] and $a=7.30$ nm [Fig. \sfref{fig:embedding}{b}], partially embedded in a glass substrate and surrounded by air, as a function of the incident wavelength under normal illumination. The degree of embedding is defined by the ratio $h/a$, with representative values taken from Fig. \sfref{fig:roots}{b)}. Three embedding scenarios were analyzed: predominantly in air ($h/a<0$, circles), mostly in glass ($h/a > 0$, diamonds), and symmetrically embedded ($h/a=0$, squares). For reference, extinction spectra of an isolated AuNP in air, computed using Mie theory, are shown as solid lines. A redshift in the resonance peak is observed with increasing embedding into the substrate, reaching up to $\sim 8$ nm for both particle sizes.

\begin{figure}[h!]\centering
	\includegraphics[scale = 1]{non-int/Fig-CextMieEmb.pdf}
	\caption{Extinction efficiencies as function of the incident wavelength of a single AuNP, partially embedded in a glass substrate and an air matrix, of radius \textbf{a)} $a = 5.78$ (red lines) nm and \textbf{b)} $a = 7.30$ nm (green lines) when illuminated at normal incidence relative to the air-glass interface. The partial embedding of the AuNPs is characterized by the quantity $h/a$ whose nominal value is specified for each AuNP according to Fig. \ref{fig:roots}. For a more complete analysis, three values of partial embedding where considered: when the AuNP is mostly embedded in air (disk), in glass (diamond) and when it is equally embedded in both media (square). }
	\label{fig:embedding}
\end{figure}

An initial numerical approach to study multiple scattering effects in a metasurface composed of AuNPs, a dimer system was modeled, consisting of two identical gold nanoparticles with radii $a = 7.30$ nm, partially embedded ($h/a=0$) at the interface between air and a glass substrate. The dimer was illuminated by a plane electromagnetic wave incident normal to the interface, with polarization aligned along the dimer’s symmetry axis. In Fig. \sfref{fig:dimmer}{a)} and \sfref{fig:dimmer}{b)} the the total electric field was calculated for the dimmer sistem considering two gap distances: 4 nm and 14 nm, respectively. Greater electric field enhancement was observed in regions adjacent to the glass substrate, the optically denser medium. This enhancement is attributed to the increased apparent area of the AuNPs within the substrate due to partial embedding. When extended to a metasurface composed of many such particles, this effect implies that the effective surface coverage is larger than what would be expected from geometry alone. Lastly, in Fig. \sfref{fig:dimmer}{c)}, the extinction cross section was computed for various interparticle gaps and compared to that of a single AuNP, calculated using Mie theory (red solid line), showing that the nearer the AuNP, the greater the redshift, reaching up to $\sim 10$ nm. When considering the contribution from the partial embedding and the multiple scattering together, spectral shifts at normal incidence can be expected to ocurre to behave as observed in the experimental extinctio spectra presented in Fig. \sfref{fig:extinction}{a)}. 


\begin{figure}[h!]\centering
	\includegraphics[scale = 1]{non-int/Fig-Dimmer.pdf}
	\caption{Magnitude of the total electric field $||\mathbf{E}_\text{Tot}||$ of a AuNP dimmer of radius $a=7.30$ nm separaeted by a gap between particles of \textbf{a)} 4 nm and \textbf{b)} 14 nm, illuminated by an incident electromagnetic plane wave traveling normal to a glass-air interface (white dashed line). The AuNP is partially embedded in the glass substrate and in the air matrix and its embedding degree is characterized by the quantity $h/a = 0$.\textbf{c)} Extinction cross section of a dimmer composed of two indentical AuNP of radius $a=7.30$ nm. The extinction cross section is presented as function of the incident wavelength of an electromagnetic plane polarized along the symmetry axis of the dimmer. The AuNPs are partially embedded in an air a matrix and a glass substrate considering $h/a = 0$, and the dimmer is illuminated at normal incidence relative to the air-glass interface. The color of each curve corresponds to different values of the in between the AuNPs of the dimmer.} 
	\label{fig:dimmer}
\end{figure}

The FEM simulations confirmed that the redshift observed in the experiments cannot be attributed solely to substrate interactions in the single-particle case, emphasizing the necessity of incorporating multiple scattering effects. Fitting with the Coherent Scattering Model (CSM) revealed an apparent polarization-dependent variation in nanoparticle size and an overall increase in surface coverage, improving agreement with experimental data. However, discrepancies between the resonance wavelengths predicted by the CSM and those measured experimentally persist, likely due to structural features not resolved in the available SEM images. Based on experimental observations, FEM simulations, and CSM predictions, it is suggested that the gold nanoparticles are partially embedded in the substrate. This partial embedding alters the apparent particle size and increases the effective surface coverage, as estimated through FEM calculations for individual particles. A comprehensive theoretical model that incorporates nanoparticle embedding into frameworks like the CSM has yet to be developed, and new approaches are still required.
